% ==============================================================================
% secoes/02_fundamentacao.tex - SEÇÃO II: FUNDAMENTAÇÃO TEÓRICA
% Tema: Federico Faggin e o Primeiro Microprocessador Monolítico (Intel 4004)
% ==============================================================================

\section{Fundamentação Teórica}
\label{sec:fundamentacao}

Conforme contextualizado na \cref{sec:introducao} e resumido nas \cref{fig:porta_mos_aluminio,fig:porta_mos_sgt}, a superação dos limites da porta metálica de alumínio exigiu reformular a geometria, a eletrostática e o projeto físico dos dispositivos PMOS \cite{cavin2012, mead1980}. Esta seção formaliza os modelos analíticos de capacitância parasita, atraso de comutação e tensão de limiar que quantificam o desempenho da tecnologia de porta autoalinhada de silício (SGT) \cite{faggin1970, sze2006}. Em seguida, examinam-se as soluções de circuito, as regras de \textit{layout} e a arquitetura de barramento multiplexado empregadas no conjunto MCS-4.

\subsection{Modelagem Física da Porta Metálica versus SGT}
\label{subsec:modelagem_fisica_sgt}

Em transistores MOS, o eletrodo de porta e o canal semicondutor atuam como placas condutoras paralelas de um capacitor separadas por uma lâmina isolante de dióxido de silício. A capacitância por unidade de área desse dielétrico ($C_{ox}$) varia na razão direta da permissividade do óxido ($\varepsilon_{ox} = \kappa_{ox}\varepsilon_0 \approx \SI{3,45e-11}{\farad\per\meter}$, produto da constante dielétrica adimensional do óxido $\kappa_{ox} \approx 3{,}9$ pela permissividade elétrica do vácuo $\varepsilon_0$) e na razão inversa da espessura física da película ($t_{ox}$):
\begin{equation}
\label{eq:cox}
C_{ox} = \frac{\varepsilon_{ox}}{t_{ox}}.
\end{equation}

A capacitância geométrica de sobreposição ($C_{ov}$) entre a porta e as difusões de fonte ou dreno decorre do produto de $C_{ox}$ pela área superposta:
\begin{equation}
\label{eq:cov}
C_{ov} = C_{ox} W L_{ov},
\end{equation}
em que $W$ indica a largura transversal do canal e $L_{ov}$ o comprimento longitudinal em que a porta recobre a região dopada adjacente. Na fabricação com porta de alumínio, em que a metalização ocorre após a formação de fonte e dreno, as tolerâncias de alinhamento das fotomáscaras acrescentam uma folga geométrica $\Delta L_{align}$ de \SIrange{3}{5}{\micro\meter} ao espalhamento térmico das impurezas:
\begin{equation}
\label{eq:lov_metal}
L_{ov,metal} = \Delta L_{align} + L_{diff,lat}.
\end{equation}
Nessa expressão, $L_{diff,lat} \approx 0{,}75 x_j$ representa a difusão lateral dos dopantes sob a borda da máscara, proporcional à profundidade vertical da junção ($x_j$) no substrato. Esse espalhamento lateral assemelha-se à tinta que avança ligeiramente sob a margem de uma fita protetora durante a aplicação. No processo SGT \cite{kerwin1969patent, faggin1970, klein1972patent}, o polissilício refratário é depositado antes das difusões. Como a própria porta atua como barreira de difusão, a folga litográfica anula-se ($\Delta L_{align} = 0$) e a sobreposição restringe-se à parcela lateral:
\begin{equation}
\label{eq:lov_sgt}
L_{ov,SGT} = L_{diff,lat}.
\end{equation}
Essa geometria (\cref{fig:modelo_capacitivo_geometria}) reduz $L_{ov}$ em três a cinco vezes e corta a capacitância parasita por borda entre 60\% e 80\% frente à porta metálica.

\begin{figure}[!t]
\centering
\includegraphics[width=\columnwidth]{modelo_capacitivo_geometria.pdf}
\caption{Geometria parametrizada da sobreposição entre porta metálica e SGT conforme as \cref{eq:cox,eq:cov,eq:lov_metal,eq:lov_sgt}, destacando $t_{ox}$, largura $W$, profundidade de junção $x_j$, difusão lateral $L_{diff,lat}$ e supressão da folga $\Delta L_{align} = 0$ na SGT.}
\label{fig:modelo_capacitivo_geometria}
\end{figure}

A redução de $C_{ov}$ atenua de forma amplificada a capacitância porta-dreno ($C_{gd}$) pelo efeito Miller durante a inversão lógica:
\begin{equation}
\label{eq:miller}
C_{in,Miller} = C_{gd} (1 - A_v) = C_{gd} (1 + |A_v|),
\end{equation}
em que $A_v = -g_m R_L$ é o ganho de tensão de pequenos sinais na região ativa, razão de amplificação entre saída e entrada cujo sinal negativo indica inversão de fase (sendo $g_m$ a transcondutância, isto é, a sensibilidade da corrente de saída em relação à tensão de porta, e $R_L$ a resistência efetiva de carga). Para módulos de ganho típicos $|A_v|$ entre 3 e 5, a capacitância física $C_{gd} \approx C_{ox} W L_{ov,d}$ (em que $L_{ov,d}$ representa a sobreposição junto ao dreno) é multiplicada por fatores de 4 a 6, penalizando o atraso de propagação ($\tau_{pd}$) da porta lógica \cite{weste2010}:
\begin{equation}
\label{eq:tau_pd}
\tau_{pd} \approx 0{,}69 R_{on} C_L \approx \frac{C_L \Delta V}{I_{on}},
\end{equation}
em que $R_{on}$ expressa a resistência equivalente do canal aberto em condução, $I_{on}$ a corrente útil fornecida para carregar ou descarregar o nó durante uma transição de amplitude $\Delta V$, e $C_L$ a capacitância total vista na saída:
\begin{equation}
\label{eq:cl_total}
C_L = C_{gd} (1 + |A_v|) + C_{gate,next} + C_{junc} + C_{wire}.
\end{equation}
Além do efeito Miller amplificado e das portas dos estágios seguintes ($C_{gate,next}$), a carga total incorpora as capacitâncias das junções $p\text{-}n$ reversamente polarizadas ($C_{junc}$, formadas na barreira entre as difusões $p^+$ e o substrato tipo $n$) e das trilhas metálicas de interconexão ($C_{wire}$). Esses dois termos operam como reservatórios secundários em paralelo que desviam parte da corrente $I_{on}$ destinada ao estágio subsequente. Ao contrair $C_{gd}$ e encurtar as ligações internas (\cref{fig:modelo_capacitivo_circuito}), o processo SGT reduziu $\tau_{pd}$ em até cinco vezes sob a mesma dissipação de potência.

\begin{figure}[!t]
\centering
\includegraphics[width=\columnwidth]{modelo_capacitivo_circuito.pdf}
\caption{Circuito equivalente de carregamento do nó de saída conforme as \cref{eq:tau_pd,eq:cl_total}, relacionando $R_{on}$, $I_{on}$ e os quatro componentes de $C_L$.}
\label{fig:modelo_capacitivo_circuito}
\end{figure}

\subsection{Eletrostática da Tensão de Limiar e Propriedades Estruturais}
\label{subsec:eletrostatica_vth_estruturais}

A tensão de limiar ($V_{th}$) de um transistor PMOS em modo de enriquecimento (dispositivo normalmente desligado quando a tensão de porta é nula) sobre silício tipo $n$ decorre do equilíbrio eletrostático exigido para estabelecer o canal condutor:
\begin{equation}
\label{eq:vth_geral}
V_{th} = V_{FB} - 2\phi_F - \frac{Q_B}{C_{ox}}.
\end{equation}
A tensão de bandas planas ($V_{FB}$) representa o potencial inicial aplicado à porta apenas para compensar diferenças internas entre materiais e cargas presas no óxido antes de iniciar a atração de lacunas, de forma semelhante a tarar uma balança antes de medir a carga útil:
\begin{equation}
\label{eq:vfb}
V_{FB} = \Phi_{ms} - \frac{Q_{ss}}{C_{ox}},
\end{equation}
sendo $\Phi_{ms} = \Phi_m - \Phi_s$ a diferença de função de trabalho entre porta e substrato, e $Q_{ss}$ a densidade superficial de cargas fixas positivas na interface óxido-silício.

O potencial de Fermi ($\phi_F$), que quantifica o afastamento energético imposto pela dopagem do substrato em relação ao silício puro, e a densidade de carga espacial na zona de depleção ($Q_B$) são definidos por:
\begin{equation}
\label{eq:phi_f}
\phi_F = \frac{k T}{q} \ln\left(\frac{N_D}{n_i}\right),
\end{equation}
\begin{equation}
\label{eq:qb}
Q_B = \sqrt{2 q \varepsilon_{Si} N_D (2\phi_F + |V_{SB}|)},
\end{equation}
em que $kT/q$ é o potencial térmico (sendo $k$ a constante de Boltzmann e $T$ a temperatura absoluta), $N_D$ a concentração volumétrica de átomos doadores de elétrons no substrato tipo $n$, $n_i$ a concentração intrínseca de portadores no silício puro, $q$ a carga elementar e $\varepsilon_{Si}$ a permissividade do silício. A grandeza $Q_B$ corresponde à carga dos íons fixos exposta na faixa do substrato que precisa ser esvaziada de elétrons livres antes da chegada das lacunas do canal. Sob polarização de corpo ($V_{SB}$, tensão reversa entre fonte e substrato), essa zona esvaziada alarga-se e reforça a barreira à condução. A substituição de \cref{eq:vfb,eq:phi_f,eq:qb} em \cref{eq:vth_geral} resulta na forma expandida de $V_{th}$ para dispositivos PMOS:
\begin{equation}
\label{eq:vth_expandido}
V_{th} = \Phi_{ms} - \frac{Q_{ss}}{C_{ox}} - 2\phi_F - \frac{\sqrt{2 q \varepsilon_{Si} N_D (2\phi_F + |V_{SB}|)}}{C_{ox}}.
\end{equation}

A substituição do alumínio por silício policristalino com dopagem $p^+$ altera diretamente $\Phi_{ms}$. Na dopagem, a concentração de boro é tão elevada que o polissilício conduz de modo semelhante a um metal. O nível energético médio dos portadores desloca-se então para o topo da banda de valência, faixa ocupada por elétrons ligados na rede cristalina. Para alumínio sobre silício tipo $n$, com $\Phi_m \approx \SI{4,10}{\electronvolt}$ e $\Phi_s \approx \SI{4,31}{\electronvolt}$, tem-se $\Phi_{ms,Al} \approx \SI{-0,21}{\volt}$. No polissilício $p^+$, a função de trabalho soma a afinidade eletrônica ($\chi \approx \SI{4,05}{\electronvolt}$, energia liberada ao capturar um elétron livre no cristal) à largura de banda proibida ($E_g \approx \SI{1,12}{\electronvolt}$, intervalo energético entre as bandas de valência e condução). Resulta assim $\Phi_{poly(p^+)} \approx \chi + E_g \approx \SI{5,17}{\electronvolt}$, elevando a diferença para $\Phi_{ms,poly} \approx \SI{5,17}{\electronvolt} - \SI{4,31}{\electronvolt} \approx \SI{+0,86}{\volt}$. A diferença líquida entre os eletrodos vale:
\begin{equation}
\label{eq:delta_phims}
\begin{aligned}
\Delta \Phi_{ms} &= \Phi_{ms,poly} - \Phi_{ms,Al} \\
&\approx \SI{+0,86}{\volt} - (\SI{-0,21}{\volt}) \approx \SI{+1,07}{\volt}.
\end{aligned}
\end{equation}
Esse ganho positivo reduz a magnitude negativa de $V_{th}$ em cerca de \SI{1,1}{\volt}, deslocando tensões típicas próximas a \SI{-3,5}{\volt} para patamares em torno de \SI{-2,4}{\volt}.

A parcela $-Q_{ss}/C_{ox}$ mantém sinal negativo em canais $p$ porque as cargas interfaciais positivas repelem as lacunas. Na orientação cristalográfica $\langle 111 \rangle$, a maior densidade superficial de ligações atômicas incompletas impõe $Q_{ss}/q \approx \SIrange{4e11}{5e11}{\per\centi\meter\squared}$. Para um óxido de \SI{100}{\nano\meter} ($C_{ox} \approx \SI{34,5}{\nano\farad\per\centi\meter\squared}$), esse termo desloca $V_{th}$ em \SI{-1,9}{\volt} a \SI{-2,3}{\volt}, elevando $|V_{th}|$ para \SIrange{3,5}{5,0}{\volt}. Na orientação $\langle 100 \rangle$, a menor densidade atômica limita as cargas a $Q_{ss}/q \approx \SIrange{0,5e11}{1,0e11}{\per\centi\meter\squared}$, restringindo o deslocamento entre \SI{-0,2}{\volt} e \SI{-0,5}{\volt}. Combinando $\Phi_{ms}$, lâminas $\langle 100 \rangle$ e recozimentos em alta temperatura, o processo SGT estabilizou $V_{th}$ entre \SI{-1,5}{\volt} e \SI{-2,5}{\volt} \cite{faggin1970, klein1972patent}, permitindo alimentar o MCS-4 com $V_{SS} = \SI{+5}{\volt}$ e $V_{DD} = \SI{-10}{\volt}$ (\SI{15}{\volt} diferenciais) em compatibilidade direta com circuitos TTL no nível alto. A \cref{tab:comparacao_eletrica} sintetiza esses parâmetros eletrostáticos, cristalográficos e térmicos.

\begin{table}[!t]
\renewcommand{\arraystretch}{1.2}
\caption{Parâmetros Eletrostáticos, Cristalográficos e Térmicos entre Porta de Alumínio e SGT}
\label{tab:comparacao_eletrica}
\centering
\begin{tabularx}{\columnwidth}{>{\raggedright\arraybackslash}X >{\centering\arraybackslash}p{2.6cm} >{\centering\arraybackslash}p{2.6cm}}
\toprule
\textbf{Parâmetro} & \textbf{Porta de Alumínio} & \textbf{Porta SGT} \\
\midrule
Diferença $\Phi_{ms}$ & $\approx \SI{-0,21}{\volt}$ & $\approx \SI{+0,86}{\volt}$ \\
Orientação do silício & $\langle 111 \rangle$ & $\langle 100 \rangle$ \\
Cargas fixas ($Q_{ss}/q$) & \SIrange{4e11}{5e11}{\per\centi\meter\squared} & \SIrange{0,5e11}{1,0e11}{\per\centi\meter\squared} \\
Deslocamento ($-Q_{ss}/C_{ox}$) & \SIrange{-1,9}{-2,3}{\volt} & \SIrange{-0,2}{-0,5}{\volt} \\
Tensão de limiar ($V_{th}$) & \SIrange{-3,5}{-5,0}{\volt} & \SIrange{-1,5}{-2,5}{\volt} \\
Limite térmico do eletrodo & Fusão \SI{660}{\celsius} e eutético \SI{577}{\celsius} & Refratário ($> \SI{1100}{\celsius}$) \\
Excursão de saída ($V_{out}$) & $V_{DD} + |V_{th}|$ (atenuada) & $V_{DD}$ (plena com \textit{bootstrapping}) \\
\bottomrule
\end{tabularx}
\end{table}

Sob a óptica de materiais, o silício policristalino produzido por deposição química em fase de vapor (CVD, técnica em que gases reagentes depositam um filme sólido uniforme sobre a lâmina aquecida) permanece estável acima de \SI{1100}{\celsius}. Essa condição permite padronizar a porta antes das difusões de fonte e dreno, conduzidas entre \SI{1050}{\celsius} e \SI{1150}{\celsius}. Em contraste, o alumínio funde a \SI{660}{\celsius} e forma com o silício uma fase eutética líquida a apenas \SI{577}{\celsius}, temperatura na qual a mistura funde abaixo do ponto de fusão de ambos os elementos puros. Acima de \SI{450}{\celsius}, o silício solubiliza-se no metal ao longo de falhas lineares na grade atômica do cristal. No resfriamento, o alumínio preenche esses vazios formando espículas condutoras (\textit{spiking}) que perfuram junções rasas ($x_j \le \SI{2}{\micro\meter}$), de modo análogo à água que rompe uma barragem fina através de fissuras, inutilizando o dispositivo por fuga reversa.

A estabilidade térmica do polissilício viabilizou ainda a captura iônica com fósforo (\textit{phosphorus gettering}) em atmosfera de oxicloreto de fósforo ($\text{POCl}_3$) entre \SI{1000}{\celsius} e \SI{1050}{\celsius}. A película gerada de vidro de fosfossilicato (PSG) atua como esponja química que fixa íons móveis de sódio ($\text{Na}^+$), evitando derivas temporais em $V_{th}$ e reduzindo fugas de corrente. Adicionalmente, o polissilício depositado sobre o óxido espesso (camada isolante mais grossa usada fora dos transistores para isolar o substrato) funciona como segunda camada condutora independente. Em circuitos de camada metálica única, o cruzamento de sinais exigia túneis difusos $p^+$ no substrato, cuja resistência de folha (resistência elétrica medida entre lados opostos de qualquer quadrado da película) situava-se entre \SI{30}{\ohm} por quadrado e \SI{100}{\ohm} por quadrado, somada a elevada capacitância de junção e perda de área. O nível intermediário de polissilício sob o alumínio eliminou esses túneis.

\subsection{Inovações de Circuito e Metodologia de \textit{Layout} no Intel 4004}
\label{subsec:inovacoes_circuito_leiaute}

Para adaptar o processo SGT às conexões irregulares da lógica aleatória do Intel 4004, Faggin combinou três recursos físicos e circuitais: os contatos enterrados, as cargas ativas com \textit{bootstrapping} dinâmico e o roteamento ortogonal em duas camadas \cite{faggin1996, faggin2021}.

Os contatos enterrados (\textit{buried contacts}), iniciados experimentalmente na Fairchild e formalizados posteriormente na patente concedida a Leslie L. Vadasz \cite{vadasz1972patent, faggin1996, faggin2021}, substituem pontes metálicas locais por ligações diretas no silício. Antes da deposição do polissilício, janelas seletivas são abertas no óxido de porta nos pontos em que o filme policristalino deve tocar o substrato monocristalino. Na difusão térmica subsequente de boro, os dopantes migram pelo polissilício para o substrato subjacente, formando uma região $p^+$ contínua que estabelece uma junção ôhmica, isto é, uma ligação resistiva simétrica sem barreira retificadora de diodo. Em \textit{layouts} convencionais de camada única (\cref{fig:contatos_enterrados}), conectar uma porta ao dreno adjacente exigia duas aberturas verticais e uma ponte de alumínio sujeita a folgas de alinhamento. A remoção dessas ilhas metálicas reduziu em até 50\% a área das células de lógica aleatória, permitindo integrar \num{2300} transistores em \SI{12}{\milli\meter\squared} \cite{faggin1996, faggin2021}.

\begin{figure}[!t]
\centering
\includegraphics[width=\columnwidth]{contatos_enterrados.pdf}
\caption{Corte transversal comparando o contato convencional com duas janelas verticais e ponte de alumínio ao contato enterrado, cuja junção ôhmica $p^+$ contínua reduz cerca de 50\% da área celular.}
\label{fig:contatos_enterrados}
\end{figure}

No domínio dos circuitos, as cargas ativas com \textit{bootstrapping} dinâmico eliminaram a perda de excursão lógica inerente a inversores PMOS convencionais \cite{faggin1972fourchip, faggin1996}. No sistema MCS-4, o potencial de repouso é $V_{SS} = \SI{+5}{\volt}$ e a alimentação ativa é negativa ($V_{DD} = \SI{-10}{\volt}$). Na configuração convencional (\cref{fig:bootstrapping_circuito}), o transistor de entrada ($M_{in}$) liga a saída ($V_{out}$) a $V_{SS}$, enquanto o transistor de carga em enriquecimento ($M_{load}$) mantém dreno e porta fixos em $V_{DD}$, atuando $V_{out}$ como sua fonte. Quando $M_{in}$ corta, $M_{load}$ descarrega $V_{out}$ de \SI{+5}{\volt} em direção a \SI{-10}{\volt} enquanto $|V_{GS}| = V_{out} - V_{DD} > |V_{th}|$. À medida que $V_{out}$ desce, a tensão porta-fonte contrai até a condução cessar em:
\begin{equation}
\label{eq:vout_max_sem_boot}
V_{out,\text{lim}} = V_{DD} + |V_{th}|.
\end{equation}
Com $V_{DD} = \SI{-10}{\volt}$ e $|V_{th}| \approx \SI{2}{\volt}$, a saída estanca próxima a \SI{-8}{\volt}, ficando aquém de $V_{DD}$ por $|V_{th}|$, o que comprime a margem de ruído e retarda a cauda da transição.

\begin{figure}[!t]
\centering
\includegraphics[width=\columnwidth]{bootstrapping_circuito.pdf}
\caption{Topologias de carga PMOS sob polaridade negativa ($V_{SS} = \SI{+5}{\volt}$ a $V_{DD} = \SI{-10}{\volt}$), comparando a carga convencional com porta fixa em $V_{DD}$ ao circuito com \textit{bootstrapping} dinâmico ($M_{iso}$, nó flutuante $N_G$, capacitor integrado $C_B$ e parasita $C_P \ll C_B$).}
\label{fig:bootstrapping_circuito}
\end{figure}

Para alcançar excursão integral até $V_{DD}$, o circuito com \textit{bootstrapping} associa um transistor de isolamento e amostragem ($M_{iso}$), a carga $M_{load}$ comandada por um nó flutuante em alta impedância ($N_G$, ponto isolado dos trilhos capaz de reter carga elétrica temporariamente) e um capacitor integrado ($C_B$) entre $V_{out}$ e $N_G$ (\cref{fig:bootstrapping_circuito}). Durante a pré-carga, etapa preparatória que fixa o estado inicial do nó, $M_{iso}$ conduz e carrega $N_G$ com uma tensão negativa $V_{G,inicial}$, entrando em corte logo após para isolar $N_G$. Quando $M_{in}$ bloqueia e $V_{out}$ sofre um degrau negativo $\Delta V_{out}$ em direção a $V_{DD}$, o capacitor $C_B$ acopla essa variação a $N_G$ por divisão capacitiva, à semelhança de uma plataforma que acompanha o deslocamento da saída, mantendo o transistor em condução até o fim do percurso:
\begin{equation}
\label{eq:delta_vg_boot}
\Delta V_G = \left( \frac{C_B}{C_B + C_P} \right) \Delta V_{out},
\end{equation}
em que $C_P$ reúne as capacitâncias parasitas em $N_G$. Dimensionando $C_B \gg C_P$, resulta $\Delta V_G \approx \Delta V_{out}$, de modo que a soma da pré-carga negativa com o degrau acoplado desloca $V_G(t)$ além de $V_{DD}$ em módulo:
\begin{equation}
\label{eq:vg_total_boot}
V_G(t) \approx V_{G,inicial} + \Delta V_{out}.
\end{equation}
Conforme mostra a \cref{fig:bootstrapping_ondas}, $V_G(t)$ mergulha abaixo de \SI{-10}{\volt}, preservando $|V_{GS}| > |V_{th}|$ durante toda a descarga. O transistor $M_{load}$ permanece na região de triodo, regime em que o canal conduz como resistor de baixa queda até igualar $V_{out} = V_{DD}$, restaurando a margem de ruído e acelerando a comutação \cite{faggin1972fourchip}.

\begin{figure}[!t]
\centering
\includegraphics[width=\columnwidth]{bootstrapping_ondas.pdf}
\caption{Formas de onda $V_{out}(t)$ e $V_G(t)$, ilustrando a estagnação em $V_{DD} + |V_{th}|$ (\SI{-8}{\volt}) sem o estágio auxiliar e a excursão plena até $V_{DD}$ (\SI{-10}{\volt}) quando o degrau negativo em $C_B$ leva $V_G(t)$ abaixo de \SI{-10}{\volt}.}
\label{fig:bootstrapping_ondas}
\end{figure}

No traçado da pastilha, cujos blocos de controle, decodificadores e somadores carecem da regularidade de memórias, Faggin estabeleceu uma disciplina de roteamento ortogonal em duas camadas sem apoio de ferramentas automáticas de síntese (\cref{fig:roteamento_ortogonal}) \cite{faggin1996, faggin2021}. As pistas de polissilício $p^+$ foram orientadas no eixo vertical e as trilhas de alumínio no eixo horizontal perpendicular. As conexões entre esses dois níveis condutores utilizavam vias verticais, janelas abertas no óxido isolante intermediário para ligar o alumínio ao polissilício, enquanto os contatos enterrados uniam o polissilício às difusões do substrato. Essa regra geométrica evitou desvios, conteve o acoplamento capacitivo entre fios e reduziu a área de silício.

\begin{figure}[!t]
\centering
\includegraphics[width=\columnwidth]{roteamento_ortogonal.pdf}
\caption{Planta superior do roteamento ortogonal em duas camadas no Intel 4004, com trilhas verticais de polissilício $p^+$ e pistas horizontais de alumínio interligadas por vias verticais e contatos enterrados sem túneis difusos.}
\label{fig:roteamento_ortogonal}
\end{figure}

\subsection{Arquitetura do Sistema MCS-4 e Barramento Multiplexado}
\label{subsec:arquitetura_mcs4_barramento}

A família MCS-4 (\textit{Micro Computer Set}) implementa um computador modular em quatro circuitos integrados fabricados em PMOS SGT de \SI{10}{\micro\meter}, acondicionados em invólucros cerâmicos DIP de 16 pinos \cite{intel1971datasheet, intel1973mcs4, noyce1981}. O processador Intel 4004 reúne cerca de \num{2300} transistores em \SI{12}{\milli\meter\squared} e executa 46 instruções sob relógio bifásico. Seu núcleo integra uma unidade lógica e aritmética (ULA, bloco que efetua somas, subtrações e operações bit a bit) de 4 bits com correção para codificação decimal binária (BCD, formato que representa dígitos de 0 a 9 em grupos de 4 bits para calculadoras), um acumulador (registrador principal que armazena resultados intermediários), um indicador de transporte (\textit{carry}, sinalizador de transbordo gerado quando uma soma excede 4 bits), 16 registradores de índice de 4 bits (memória interna rápida endereçável em pares de 8 bits) e um contador de programa de 12 bits com pilha de três níveis (apontador de endereço da próxima instrução com três posições de retorno para sub-rotinas).

Os demais componentes fornecem memória e expansão de entrada e saída sobre o mesmo barramento (\cref{fig:mcs4_arquitetura_blocos}). O Intel 4001 incorpora \num{2048} bits de ROM mascarada (memória somente de leitura gravada permanentemente na fábrica durante a fotolitografia), dispostos em 256 palavras de 8 bits e uma porta de entrada e saída de 4 bits, admitindo até 16 unidades no sistema (4 quilobytes). O Intel 4002 provê \num{320} bits de RAM dinâmica (memória de leitura e escrita que retém dados em capacitores internos renovados periodicamente), distribuídos em quatro registradores de 20 caracteres de 4 bits com porta de saída e seleção por quatro linhas CM-RAM. O Intel 4003 atua como expansor estático de saída (circuito que amplia as linhas disponíveis mantendo os dados retidos) mediante um registrador de deslocamento de 10 bits, cadeia que recebe bits em série por um único fio e os disponibiliza simultaneamente em saídas paralelas para periféricos.

Em 1971, limites de custo e rendimento confinavam o encapsulamento comercial a 16 pinos, inviabilizando barramentos distintos para endereços, instruções e dados. Para contornar essa restrição, adotou-se um barramento bidirecional multiplexado no tempo de 4 bits ($D_0$ a $D_3$) \cite{intel1971datasheet, intel1973mcs4, faggin1974patent}. Os 16 terminais do Intel 4004 compreendem quatro pinos de dados ($D_0$ a $D_3$), duas entradas de relógio ($\phi_1$ e $\phi_2$), saída de sincronização (\textit{SYNC}), entradas de reinicialização (\textit{RESET}) e teste lógico (\textit{TEST}), quatro linhas de seleção de RAM (CM-RAM0 a CM-RAM3), habilitação de ROM (CM-ROM) e dois terminais de alimentação ($V_{SS} = \SI{+5}{\volt}$ e $V_{DD} = \SI{-10}{\volt}$).

\begin{figure}[!t]
\centering
\includegraphics[width=\columnwidth]{mcs4_arquitetura_blocos.pdf}
\caption{Diagrama em blocos dos circuitos integrados 4004, 4001, 4002 e 4003 do sistema MCS-4, exibindo a pinagem de 16 pinos da CPU, o barramento bidirecional de 4 bits ($D_0$ a $D_3$), as linhas de controle e os trilhos $V_{SS} = \SI{+5}{\volt}$ e $V_{DD} = \SI{-10}{\volt}$.}
\label{fig:mcs4_arquitetura_blocos}
\end{figure}

O sincronismo baseia-se em um relógio bifásico ($\phi_1$ e $\phi_2$), composto por dois pulsos alternados que nunca permanecem ativos juntos, à semelhança das comportas de uma eclusa em que uma fecha antes de a outra abrir para evitar conflitos elétricos no barramento. Com período nominal de \SI{1,35}{\micro\second} (\SI{740,7}{\kilo\hertz}), uma instrução de uma palavra (8 bits) requer oito períodos de relógio, definindo um ciclo de \SI{10,8}{\micro\second} (aproximadamente \num{92600} instruções por segundo), enquanto instruções de duas palavras (16 bits) consomem \SI{21,6}{\micro\second} \cite{intel1971datasheet, intel1973mcs4}. Conforme ilustra a \cref{fig:mcs4_cronograma_subciclos}, o ciclo de \SI{10,8}{\micro\second} reparte-se em oito subciclos:
\begin{enumerate}
\item A1 (\textit{Address 1}): envio dos quatro bits menos significativos de endereço da ROM ($A_0$ a $A_3$).
\item A2 (\textit{Address 2}): envio dos quatro bits intermediários de endereço da ROM ($A_4$ a $A_7$).
\item A3 (\textit{Address 3}): envio dos quatro bits mais significativos de endereço da ROM ($A_8$ a $A_{11}$).
\item M1 (\textit{Memory 1}): leitura dos quatro bits superiores da instrução na ROM, código de operação (OPR).
\item M2 (\textit{Memory 2}): leitura dos quatro bits inferiores da instrução na ROM, modificador de operação (OPA).
\item X1 (\textit{Execute 1}): decodificação da instrução e início da execução lógica ou aritmética na ULA e registradores.
\item X2 (\textit{Execute 2}): conclusão das operações na ULA, atualização do acumulador e do indicador \textit{carry}, além do envio de endereço ou dado em operações com RAM e portas.
\item X3 (\textit{Execute 3}): conclusão do ciclo de execução, eventuais transferências de entrada ou saída com a RAM 4002 e emissão do pulso \textit{SYNC} para anunciar o ciclo seguinte.
\end{enumerate}

\begin{figure}[!t]
\centering
\includegraphics[width=\columnwidth]{mcs4_cronograma_subciclos.pdf}
\caption{Cronograma de \SI{10,8}{\micro\second} do sistema MCS-4 sob relógio bifásico não sobreposto ($\phi_1$, $\phi_2$ a \SI{740,7}{\kilo\hertz}), sinal $\text{SYNC}$ e os oito subciclos A1 a X3 divididos em duas faixas.}
\label{fig:mcs4_cronograma_subciclos}
\end{figure}
